\documentclass{article}%
\usepackage[T1]{fontenc}%
\usepackage[utf8]{inputenc}%
\usepackage{lmodern}%
\usepackage{textcomp}%
\usepackage{lastpage}%
\usepackage{geometry}%
\geometry{margin=1in,headheight=14pt,headsep=25pt}%
\usepackage{hyperref}%
\usepackage{enumitem}%
\usepackage{fancyhdr}%
\usepackage{xcolor}%
\usepackage{url}%
\usepackage{breakurl}%
%

            \hypersetup{
                colorlinks=true,
                linkcolor=blue,
                filecolor=magenta,
                urlcolor=blue
            }
            \pagestyle{fancy}
            \fancyhf{}
            \rhead{Generated on \today}
            \cfoot{\thepage}
            
            % Configure enumeration settings
            \setlist[enumerate,1]{label=\arabic*.}
            \setlist[enumerate,2]{label=\alph*.}
            \setlist[enumerate,3]{label=\Alph*.}
            \setlist[enumerate]{itemsep=0.5em}
            
            % Define a command for red text
            \newcommand{\incorrect}[1]{\textcolor{red}{#1}}
        %
\lhead{2024{-}10{-}28{-}Farkas}%
\title{Practice Questions for 2024{-}10{-}28{-}Farkas}%
\author{Generated by Scribe.AI}%
\date{\today}%
%
\begin{document}%
\normalsize%
\maketitle%
\section{Questions}%
\label{sec:Questions}%
\begin{enumerate}%
\item Consider Farkas' Lemma as presented in Slide 1.  Suppose we have a system of inequalities $Ax \le b$ where $A$ is a $3 \times 2$ matrix and $b$ is a $3 \times 1$ vector.  If this system has no solution, what can we definitively conclude about the dimensions of a vector $y$ satisfying the conditions $A^T y = 0$, $y \ge 0$, and $b^T y < 0$, and why?%
\begin{enumerate}%
\item $y$ is a $2 \times 1$ vector.%
\item $y$ is a $3 \times 1$ vector.%
\item $y$ is a $2 \times 3$ vector.%
\item $y$ is a $3 \times 2$ vector.%
\item The dimensions of $y$ cannot be determined from the given information.%
\end{enumerate}%
\vspace{1em}%
\item%
\begin{enumerate}%
\item Consider the system of inequalities $Ax \le b$, where $A$ is a $m \times n$ matrix and $b$ is a $m \times 1$ vector.  According to Farkas' Lemma (Slide 1), if this system has no solution, what can we conclude about the existence of a vector $y$ satisfying $A^T y = 0$, $y \ge 0$, and $b^T y < 0$?%
\begin{enumerate}%
\item Such a vector $y$ always exists.%
\item Such a vector $y$ may or may not exist.%
\item Such a vector $y$ never exists.%
\item The existence of $y$ depends on the specific values of $A$ and $b$.%
\item The existence of $y$ depends on whether $m > n$ or $m < n$.%
\end{enumerate}%
\vspace{0.5em}%
\item Suppose we have a system $Ax = b$ where $A$ is a $5 \times 3$ matrix and $b$ is a $5 \times 1$ vector.  According to one version of Farkas' Lemma (Slide 2), if there is no non-negative solution $x$ (i.e., $x \ge 0$), what can we conclude about the existence of a vector $y$ such that $y^T A \ge 0^T$ and $y^T b < 0$?%
\begin{enumerate}%
\item Such a vector $y$ always exists.%
\item Such a vector $y$ may or may not exist.%
\item Such a vector $y$ never exists.%
\item The existence of $y$ depends on the specific values of $A$ and $b$.%
\item The existence of $y$ depends on the rank of matrix $A$.%
\end{enumerate}%
\vspace{0.5em}%
\item Consider the three variants of Farkas' Lemma presented in Slide 4.  Which variant directly addresses the solvability of the system $Ax \le b$ without any non-negativity constraints on $x$?%
\begin{enumerate}%
\item Variant (i)%
\item Variant (ii)%
\item Variant (iii)%
\item Variants (i) and (ii)%
\item None of the variants directly address this.%
\end{enumerate}%
\vspace{0.5em}%
\end{enumerate}%
\item Consider Farkas' Lemma (Slide 1).  If the system $Ax \le b$ has no solution, where $A$ is a $m \times n$ matrix and $b$ is a $m \times 1$ vector, what can we conclude about the vector $y$ satisfying $A^T y = 0$, $y \ge 0$, and $b^T y < 0$?  Specifically, what is the dimension of $y$?%
\begin{enumerate}%
\item The dimension of $y$ is $n \times 1$.%
\item The dimension of $y$ is $n \times m$.%
\item The dimension of $y$ is $m \times 1$.%
\item The dimension of $y$ is $m \times n$.%
\item The dimension of $y$ cannot be determined from the given information.%
\end{enumerate}%
\vspace{1em}%
\item Consider the system $Ax = b$, where $A$ is a $5 \times 3$ matrix and $b$ is a $5 \times 1$ vector (Slide 2).  According to Farkas' Lemma, if there is no non-negative solution $x$ ($x \ge 0$), what can we conclude about the vector $y$ such that $y^T A \ge 0^T$ and $y^T b < 0$? What is the dimension of $y$?%
\begin{enumerate}%
\item The dimension of $y$ is $3 \times 1$.%
\item The dimension of $y$ is $3 \times 5$.%
\item The dimension of $y$ is $5 \times 1$.%
\item The dimension of $y$ is $5 \times 3$.%
\item The dimension of $y$ cannot be determined.%
\end{enumerate}%
\vspace{1em}%
\item Referring to Slide 4, which variant of Farkas' Lemma directly addresses the solvability of the system $Ax \le b$ without any non-negativity constraints on $x$?%
\begin{enumerate}%
\item Variant (i)%
\item Variant (ii)%
\item Variant (iii)%
\item None of the variants address this directly.%
\item Variants (i) and (ii) both address this.%
\end{enumerate}%
\vspace{1em}%
\item Suppose we have a system of inequalities $Ax \le b$, where $A$ is a $3 \times 2$ matrix and $b$ is a $3 \times 1$ vector.  If this system has no solution, what are the dimensions of the vector $y$ that satisfies $A^T y = 0$, $y \ge 0$, and $b^T y < 0$ (Slide 1), and what is the interpretation of this result in the context of Farkas' Lemma?%
\begin{enumerate}%
\item $y$ is $2 \times 1$; it indicates the system is infeasible.%
\item $y$ is $3 \times 1$; it indicates the system is infeasible.%
\item $y$ is $2 \times 1$; it indicates the system is feasible.%
\item $y$ is $3 \times 1$; it indicates the system is feasible.%
\item The dimensions of $y$ cannot be determined.%
\end{enumerate}%
\vspace{1em}%
\item%
\begin{enumerate}%
\item Consider Farkas' Lemma (Slide 1): The system $Ax \le b$ has no solutions if and only if there exists a vector $y$ such that $A^T y = 0$, $y \ge 0$, and $b^T y < 0$.  Suppose $A$ is a $m \times n$ matrix and $b$ is a $m \times 1$ vector. What is the dimension of the vector $y$, and why?%
\begin{enumerate}%
\item $n \times 1$%
\item $m \times 1$%
\item $n \times m$%
\item $m \times n$%
\item The dimension of $y$ is not determined by the dimensions of $A$ and $b$.%
\end{enumerate}%
\vspace{0.5em}%
\item Consider the alternative version of Farkas' Lemma in Slide 2:  Exactly one of the following holds: (F1) There exists $x \in \mathbb{R}^n$ such that $Ax = b$ and $x \ge 0$; (F2) There exists $y \in \mathbb{R}^m$ such that $y^T A \ge 0^T$ and $y^T b < 0$. If $A$ is a $5 \times 3$ matrix and $b$ is a $5 \times 1$ vector, what are the dimensions of $x$ and $y$, and why?%
\begin{enumerate}%
\item $x$ is $3 \times 1$, $y$ is $3 \times 1$%
\item $x$ is $5 \times 1$, $y$ is $5 \times 1$%
\item $x$ is $3 \times 1$, $y$ is $5 \times 1$%
\item $x$ is $5 \times 1$, $y$ is $3 \times 1$%
\item $x$ is $3 \times 5$, $y$ is $5 \times 3$%
\end{enumerate}%
\vspace{0.5em}%
\item Referring to Slide 4, which variant of Farkas' Lemma directly addresses the solvability of the system $Ax \le b$ without any non-negativity constraints on $x$?%
\begin{enumerate}%
\item Variant (i)%
\item Variant (ii)%
\item Variant (iii)%
\item Variants (i) and (ii)%
\item None of the variants directly address this.%
\end{enumerate}%
\vspace{0.5em}%
\end{enumerate}%
\end{enumerate}

%
\newpage%
\section{Answers}%
\label{sec:Answers}%
\begin{enumerate}%
\item Consider Farkas' Lemma as presented in Slide 1.  Suppose we have a system of inequalities $Ax \le b$ where $A$ is a $3 \times 2$ matrix and $b$ is a $3 \times 1$ vector.  If this system has no solution, what can we definitively conclude about the dimensions of a vector $y$ satisfying the conditions $A^T y = 0$, $y \ge 0$, and $b^T y < 0$, and why?%
\begin{enumerate}%
\item \incorrect{Incorrect.  The dimensions of $y$ are determined by the dimensions of $A^T$, which is the transpose of $A$.}%
\item The correct answer is B.  Since $A$ is a $3 \times 2$ matrix, its transpose $A^T$ is a $2 \times 3$ matrix.  The equation $A^T y = 0$ implies that $y$ must have 3 rows to allow for matrix multiplication. The condition $y \ge 0$ indicates that $y$ is a column vector, thus $y$ is a $3 \times 1$ vector.%
\item \incorrect{Incorrect. $y$ must be a column vector to satisfy the condition $A^T y = 0$.}%
\item \incorrect{Incorrect.  The dimensions of $y$ are determined by the dimensions of $A^T$, which is the transpose of $A$.}%
\item \incorrect{Incorrect. The dimensions of $y$ are directly determined by the dimensions of $A$ and the requirement that $A^T y = 0$ be a valid matrix multiplication.}%
\end{enumerate}%
\vspace{1em}%
\item%
\begin{enumerate}%
\item Consider the system of inequalities $Ax \le b$, where $A$ is a $m \times n$ matrix and $b$ is a $m \times 1$ vector.  According to Farkas' Lemma (Slide 1), if this system has no solution, what can we conclude about the existence of a vector $y$ satisfying $A^T y = 0$, $y \ge 0$, and $b^T y < 0$?%
\begin{enumerate}%
\item Farkas' Lemma states that if the system $Ax \le b$ has no solution, then there exists a vector $y$ such that $A^T y = 0$, $y \ge 0$, and $b^T y < 0$. This is a direct consequence of the lemma's statement.%
\item \incorrect{This is incorrect because Farkas' Lemma provides a definitive statement about the existence of $y$ when $Ax \le b$ has no solution.}%
\item \incorrect{This is incorrect; Farkas' Lemma guarantees the existence of such a vector $y$ when the original system has no solution.}%
\item \incorrect{This is incorrect. The lemma's statement is independent of the specific values of $A$ and $b$, provided the system $Ax \le b$ is infeasible.}%
\item \incorrect{This is incorrect. The dimensions of $A$ and $b$ influence the dimensions of $y$, but the lemma's core statement holds regardless of the relationship between $m$ and $n$.}%
\end{enumerate}%
\vspace{0.5em}%
\item Suppose we have a system $Ax = b$ where $A$ is a $5 \times 3$ matrix and $b$ is a $5 \times 1$ vector.  According to one version of Farkas' Lemma (Slide 2), if there is no non-negative solution $x$ (i.e., $x \ge 0$), what can we conclude about the existence of a vector $y$ such that $y^T A \ge 0^T$ and $y^T b < 0$?%
\begin{enumerate}%
\item According to the version of Farkas' Lemma on Slide 2, if the system $Ax = b$ has no non-negative solution, then there exists a vector $y$ such that $y^T A \ge 0^T$ and $y^T b < 0$. This is a direct application of the lemma.%
\item \incorrect{This is incorrect; Farkas' Lemma provides a definitive statement about the existence of $y$ when $Ax = b$ has no non-negative solution.}%
\item \incorrect{This is incorrect; the lemma guarantees the existence of such a vector $y$ when the original system has no non-negative solution.}%
\item \incorrect{This is incorrect. The lemma's statement is independent of the specific values of $A$ and $b$, provided the system $Ax = b$ has no non-negative solution.}%
\item \incorrect{While the rank of $A$ is relevant to the solvability of $Ax = b$, the lemma's statement about the existence of $y$ is independent of the rank of $A$.}%
\end{enumerate}%
\vspace{0.5em}%
\item Consider the three variants of Farkas' Lemma presented in Slide 4.  Which variant directly addresses the solvability of the system $Ax \le b$ without any non-negativity constraints on $x$?%
\begin{enumerate}%
\item \incorrect{Variant (i) deals with $Ax = b$ and $x \ge 0$.}%
\item \incorrect{Variant (ii) deals with $Ax \le b$ and $x \ge 0$.}%
\item Variant (iii) explicitly deals with the existence of a solution to $Ax \le b$ without any restriction on the sign of $x$.%
\item \incorrect{Both (i) and (ii) impose non-negativity constraints on $x$.}%
\item \incorrect{This is incorrect; variant (iii) directly addresses this case.}%
\end{enumerate}%
\vspace{0.5em}%
\end{enumerate}%
\item Consider Farkas' Lemma (Slide 1).  If the system $Ax \le b$ has no solution, where $A$ is a $m \times n$ matrix and $b$ is a $m \times 1$ vector, what can we conclude about the vector $y$ satisfying $A^T y = 0$, $y \ge 0$, and $b^T y < 0$?  Specifically, what is the dimension of $y$?%
\begin{enumerate}%
\item \incorrect{Incorrect. The dimension of $y$ is determined by the dimensions of $A^T$, which is $n \times m$.  The equation $A^T y = 0$ requires $y$ to have $m$ entries.}%
\item \incorrect{Incorrect.  This is not a valid matrix dimension for the equation $A^T y = 0$ to hold.}%
\item The correct answer is C.  Farkas' Lemma states that if $Ax \le b$ has no solution, then there exists a vector $y$ such that $A^T y = 0$, $y \ge 0$, and $b^T y < 0$. Since $A$ is an $m \times n$ matrix, $A^T$ is an $n \times m$ matrix.  For the equation $A^T y = 0$ to be consistent, $y$ must have dimension $m \times 1$.%
\item \incorrect{Incorrect. This is not a valid matrix dimension for the equation $A^T y = 0$ to hold.}%
\item \incorrect{Incorrect. The dimension of $y$ is directly determined by the dimensions of $A$ and the requirement that $A^T y = 0$ be a valid matrix equation.}%
\end{enumerate}%
\vspace{1em}%
\item Consider the system $Ax = b$, where $A$ is a $5 \times 3$ matrix and $b$ is a $5 \times 1$ vector (Slide 2).  According to Farkas' Lemma, if there is no non-negative solution $x$ ($x \ge 0$), what can we conclude about the vector $y$ such that $y^T A \ge 0^T$ and $y^T b < 0$? What is the dimension of $y$?%
\begin{enumerate}%
\item \incorrect{Incorrect.  $y$ must have 5 entries to make $y^T A$ a valid matrix multiplication.}%
\item \incorrect{Incorrect. This is not a valid matrix dimension for the product $y^T A$ to be defined.}%
\item The correct answer is C.  Farkas' Lemma (Slide 2) states that if $Ax = b$ has no non-negative solution, then there exists a vector $y$ such that $y^T A \ge 0^T$ and $y^T b < 0$. Since $A$ is a $5 \times 3$ matrix, $y^T$ must be a $1 \times 5$ row vector for the product $y^T A$ to be defined. Therefore, $y$ is a $5 \times 1$ column vector.%
\item \incorrect{Incorrect. This is not a valid matrix dimension for the product $y^T A$ to be defined.}%
\item \incorrect{Incorrect. The dimension of $y$ is determined by the requirement that $y^T A$ be a valid matrix product.}%
\end{enumerate}%
\vspace{1em}%
\item Referring to Slide 4, which variant of Farkas' Lemma directly addresses the solvability of the system $Ax \le b$ without any non-negativity constraints on $x$?%
\begin{enumerate}%
\item \incorrect{Incorrect. Variant (i) deals with $Ax = b$ and non-negative $x$.}%
\item \incorrect{Incorrect. Variant (ii) deals with $Ax \le b$ and non-negative $x$.}%
\item The correct answer is C. Variant (iii) of Farkas' Lemma (Slide 4) states that the system $Ax \le b$ has a solution (without any non-negativity constraint on $x$) if and only if every non-negative $y \in \mathbb{R}^m$ with $y^T A = 0^T$ also satisfies $y^T b \ge 0$.%
\item \incorrect{Incorrect. Variant (iii) directly addresses this case.}%
\item \incorrect{Incorrect. Only variant (iii) addresses the solvability of $Ax \le b$ without non-negativity constraints on $x$.}%
\end{enumerate}%
\vspace{1em}%
\item Suppose we have a system of inequalities $Ax \le b$, where $A$ is a $3 \times 2$ matrix and $b$ is a $3 \times 1$ vector.  If this system has no solution, what are the dimensions of the vector $y$ that satisfies $A^T y = 0$, $y \ge 0$, and $b^T y < 0$ (Slide 1), and what is the interpretation of this result in the context of Farkas' Lemma?%
\begin{enumerate}%
\item \incorrect{Incorrect. The dimension of $y$ must match the number of rows in $A$, which is 3.}%
\item The correct answer is B.  Since $A$ is a $3 \times 2$ matrix, $A^T$ is a $2 \times 3$ matrix.  For the equation $A^T y = 0$ to be consistent, $y$ must be a $3 \times 1$ vector.  Farkas' Lemma states that if $Ax \le b$ has no solution, then such a $y$ exists, confirming the infeasibility of the original system.%
\item \incorrect{Incorrect. The existence of such a $y$ indicates infeasibility, not feasibility.}%
\item \incorrect{Incorrect. While the dimension is correct, the interpretation is wrong; the existence of $y$ signifies infeasibility.}%
\item \incorrect{Incorrect. The dimensions of $y$ are directly determined by the dimensions of $A$.}%
\end{enumerate}%
\vspace{1em}%
\item%
\begin{enumerate}%
\item Consider Farkas' Lemma (Slide 1): The system $Ax \le b$ has no solutions if and only if there exists a vector $y$ such that $A^T y = 0$, $y \ge 0$, and $b^T y < 0$.  Suppose $A$ is a $m \times n$ matrix and $b$ is a $m \times 1$ vector. What is the dimension of the vector $y$, and why?%
\begin{enumerate}%
\item \incorrect{Incorrect.  The transpose of A is $n \times m$, so $y$ cannot be $n \times 1$ for the multiplication $A^T y$ to be defined.}%
\item The dimension of $y$ is $m \times 1$.  The equation $A^T y = 0$ involves the transpose of $A$, which is an $n \times m$ matrix.  Since $y$ is multiplied by $A^T$ on the left, $y$ must be a column vector with $m$ rows to allow for matrix multiplication.  Therefore, the dimension of $y$ is $m \times 1$.%
\item \incorrect{Incorrect.  $A^T$ is $n \times m$, so $y$ cannot be $n \times m$ for the multiplication $A^T y$ to be defined.}%
\item \incorrect{Incorrect.  $A^T$ is $n \times m$, so $y$ cannot be $m \times n$ for the multiplication $A^T y$ to be defined.}%
\item \incorrect{Incorrect. The dimension of $y$ is directly determined by the dimensions of $A$ and the requirement that $A^T y$ is a well-defined matrix multiplication.}%
\end{enumerate}%
\vspace{0.5em}%
\item Consider the alternative version of Farkas' Lemma in Slide 2:  Exactly one of the following holds: (F1) There exists $x \in \mathbb{R}^n$ such that $Ax = b$ and $x \ge 0$; (F2) There exists $y \in \mathbb{R}^m$ such that $y^T A \ge 0^T$ and $y^T b < 0$. If $A$ is a $5 \times 3$ matrix and $b$ is a $5 \times 1$ vector, what are the dimensions of $x$ and $y$, and why?%
\begin{enumerate}%
\item \incorrect{Incorrect. The dimensions of $x$ and $y$ are not both $3 \times 1$.}%
\item \incorrect{Incorrect. The dimension of $x$ is not $5 \times 1$.}%
\item In (F1), $Ax = b$ requires $A$ ($5 \times 3$) multiplied by $x$ to result in a $5 \times 1$ vector $b$. Thus, $x$ must be $3 \times 1$. In (F2), $y^T A$ must be a $1 \times 3$ row vector, and $y^T b$ must be a scalar. Since $A$ has 5 rows, $y$ must be a $5 \times 1$ column vector.%
\item \incorrect{Incorrect. The dimension of $y$ is not $3 \times 1$.}%
\item \incorrect{Incorrect. The dimensions of $x$ and $y$ are not $3 \times 5$ and $5 \times 3$ respectively.}%
\end{enumerate}%
\vspace{0.5em}%
\item Referring to Slide 4, which variant of Farkas' Lemma directly addresses the solvability of the system $Ax \le b$ without any non-negativity constraints on $x$?%
\begin{enumerate}%
\item \incorrect{Incorrect. Variant (i) deals with $Ax = b$ and $x \ge 0$.}%
\item \incorrect{Incorrect. Variant (ii) deals with $Ax \le b$ and $x \ge 0$.}%
\item Variant (iii) states: "The system $Ax \le b$ has a solution if and only if every nonnegative $y \in \mathbb{R}^m$ with $y^T A = 0^T$ also satisfies $y^T b \ge 0$."  This is the only variant that doesn't impose a non-negativity constraint on $x$.%
\item \incorrect{Incorrect. Only variant (iii) addresses the solvability of $Ax \le b$ without non-negativity constraints on $x$.}%
\item \incorrect{Incorrect. Variant (iii) directly addresses this.}%
\end{enumerate}%
\vspace{0.5em}%
\end{enumerate}%
\end{enumerate}

%
\end{document}